\subsection{第二检测点候选区域}

记第一次检测点为 \(S_1\)，其示向中心射线方向为 \(\vect u_1\)。由一次测向形成的源位置可行域为
\begin{equation}
  \Omega_1=W_1\cap B.
\end{equation}
第二检测点 \(S_2\) 首先应位于目标圆域内。考虑有效接收半径未知但位于 \([1000,1500]\,\mathrm{m}\)，定义可能接收域
\begin{equation}
  \mathcal{R}_{\mathrm{poss}}
  =\left\{s\in B:\min_{g\in\Omega_1}\|s-g\|_2\le1500\right\},
\end{equation}
以及保守接收域
\begin{equation}
  \mathcal{R}_{\mathrm{safe}}
  =\left\{s\in B:\max_{g\in\Omega_1}\|s-g\|_2\le1000\right\}.
\end{equation}

为避免基线过短，并鼓励接近垂直交会，给定基线区间 \([L_{\min},L_{\max}]\) 和允许的偏离角 \(\delta\)，候选区域可写为
\begin{equation}
  \mathcal{C}=\left\{s\in\mathcal{R}_{\mathrm{poss}}:
  L_{\min}\le\|s-S_1\|_2\le L_{\max},\quad
  |(s-S_1)^{\mathsf T}\vect u_1|
  \le\|s-S_1\|_2\sin\delta\right\}.
  \label{eq:candidate-region}
\end{equation}
若 \(\mathcal{R}_{\mathrm{safe}}\cap\mathcal{C}\neq\varnothing\)，优先在该交集内选点；否则在 \(\mathcal{C}\) 中以接收风险作为惩罚项优化。
